% Options for packages loaded elsewhere
\PassOptionsToPackage{unicode}{hyperref}
\PassOptionsToPackage{hyphens}{url}
\documentclass[
  11pt,
]{article}
\usepackage{xcolor}
\usepackage[margin=1in]{geometry}
\usepackage{amsmath,amssymb}
\setcounter{secnumdepth}{3}
\usepackage{iftex}
\ifPDFTeX
  \usepackage[T1]{fontenc}
  \usepackage[utf8]{inputenc}
  \usepackage{textcomp} % provide euro and other symbols
\else % if luatex or xetex
  \usepackage{unicode-math} % this also loads fontspec
  \defaultfontfeatures{Scale=MatchLowercase}
  \defaultfontfeatures[\rmfamily]{Ligatures=TeX,Scale=1}
\fi
\usepackage{lmodern}
\ifPDFTeX\else
  % xetex/luatex font selection
\fi
% Use upquote if available, for straight quotes in verbatim environments
\IfFileExists{upquote.sty}{\usepackage{upquote}}{}
\IfFileExists{microtype.sty}{% use microtype if available
  \usepackage[]{microtype}
  \UseMicrotypeSet[protrusion]{basicmath} % disable protrusion for tt fonts
}{}
\makeatletter
\@ifundefined{KOMAClassName}{% if non-KOMA class
  \IfFileExists{parskip.sty}{%
    \usepackage{parskip}
  }{% else
    \setlength{\parindent}{0pt}
    \setlength{\parskip}{6pt plus 2pt minus 1pt}}
}{% if KOMA class
  \KOMAoptions{parskip=half}}
\makeatother
\usepackage{longtable,booktabs,array}
\newcounter{none} % for unnumbered tables
\usepackage{calc} % for calculating minipage widths
% Correct order of tables after \paragraph or \subparagraph
\usepackage{etoolbox}
\makeatletter
\patchcmd\longtable{\par}{\if@noskipsec\mbox{}\fi\par}{}{}
\makeatother
% Allow footnotes in longtable head/foot
\IfFileExists{footnotehyper.sty}{\usepackage{footnotehyper}}{\usepackage{footnote}}
\makesavenoteenv{longtable}
\setlength{\emergencystretch}{3em} % prevent overfull lines
\providecommand{\tightlist}{%
  \setlength{\itemsep}{0pt}\setlength{\parskip}{0pt}}
\usepackage{bookmark}
\IfFileExists{xurl.sty}{\usepackage{xurl}}{} % add URL line breaks if available
\urlstyle{same}
\hypersetup{
  hidelinks,
  pdfcreator={LaTeX via pandoc}
}

\newtheorem{theorem}{Theorem}
\newtheorem{lemma}{Lemma}
\newtheorem{proposition}{Proposition}
\newtheorem{corollary}{Corollary}
\newtheorem{definition}{Definition}
\newtheorem{remark}{Remark}

\author{}
\date{}

\begin{document}

\title{From a Single Axiom to Four-Dimensional Spacetime: A Readout Mechanism for Discrete Ontology}
\author{}
\date{}
\maketitle

\begin{abstract}

This work establishes a discrete ontological framework whose sole
starting point is a single axiom: \emph{zero never stays still}. No
manifold, metric, probability, or action is introduced. Words, graphs,
divergence identities, zero-sum conservation, and integer multiplicities
are all derived structures. We establish a layered index system
distinguishing the evolution layer from the historical layer, and give a
strict parameter account: one axiom, one free parameter, two named
inputs. On this basis we present the main derivation chain: zero-sum
constraint \(\to\) space of balanced words \(\to\) self-identifying
coarse-graining partition \(\to\) spectral dimension readout. At
\(n=16\), the participation index of the height-level partition is
\(q=\sum_b\omega_b^2=0.55777859\), which lies in the interval
\((1/2,3/5)\); the functional \(F_D=\binom{D}{2}q^D\) then has a unique
finite maximum at \(D=4\). We explicitly list all open items and
conditional boundaries: the necessity of the quantum contextuality
threshold is unproven; the partition criterion has been shown
non-unique; and \(D=4\) holds only for \(n\in\{16,20\}\).

\end{abstract}

\noindent\textbf{Keywords}: discrete ontology; zero-sum constraint;
coarse-graining; spacetime dimension; spectral dimension; KCBS
inequality

\begin{center}\rule{0.5\linewidth}{0.5pt}\end{center}

\section{Introduction}\label{introduction}

The spacetime dimension \(D=4\) is an input rather than a conclusion
within the standard physical framework {[}1,2{]}. The axiomatic systems
of quantum mechanics and general relativity likewise lack a common
derivation basis. Our approach is to start from a discrete substrate
containing no geometry, no probability, and no action, such that
dimension becomes a computable readout quantity rather than an intrinsic
property of microscopic degrees of freedom.

The core stance is: dimension in discrete systems is inherently a
readout quantity at the coarse-graining scale, not a microscopic degree
of freedom. The framework therefore does not posit a ``microscopic
dimension''; it only posits a ``readout dimension.'' The method proceeds
in four steps: axiom \(\to\) combinatorial structure \(\to\) partition
\(\to\) readout.

The paper is organized as follows: §2 states the single axiom and its
mathematical formulation; §3 derives words, graphs, the divergence
dictionary, and zero-sum conservation; §4 establishes layered indices;
§5 gives the object dictionary; §6 registers the parameter account; §7
presents the main derivation chain from balanced words to dimension
readout; §8 summarizes proven results; §9 lists conditions and open
items; §10 is discussion.

\begin{center}\rule{0.5\linewidth}{0.5pt}\end{center}

\section{The Axiom and Its Mathematical
Formulation}\label{the-axiom-and-its-mathematical-formulation}

\subsection{The Single Axiom}\label{the-single-axiom}

\textbf{Axiom 1 (Zero Never Stays Still).} The equilibrium (closed)
state is not stationary; it moves continuously. Motion does not
terminate at any point in time. Motion carries no weight---there is no
selection rule in the system.

\subsection{Mathematical
Formulation}\label{mathematical-formulation}

The three postulates of the axiom correspond to the following
structures:

\[
Z_0 = \underbrace{(\{\pm 1\},\,\iota)}_{\text{sign structure}} + \underbrace{(\mathcal{W}_n,\,\mu)}_{\text{closure constraint + counting measure}}.
\tag{2.1}
\]

\textbf{Sign structure.} Motion has two directions, giving
\(\{\pm 1\}\). Zero is not a stationary point but an equilibrium
relation, and is therefore expressed as the quotient by an involution
\(\iota:\{\pm1\}\to\{\pm1\}\) (\(\iota^2=\mathrm{id}\),
\(\iota(x)\neq x\)):

\[
\text{zero} = \{\pm1\}/\iota.
\tag{2.2}
\]

\textbf{Closure constraint.} Motion along a graph \(\Gamma=(C,E)\) gives
a divergence \(d_w(v)=\#\{\text{entering }v\}-\#\{\text{leaving }v\}\)
satisfying the zero-sum identity

\[
\sum_{v\in C} d_w(v) \equiv 0, \qquad \text{closed} \iff d_w = 0.
\tag{2.3}
\]

\textbf{Counting measure.} No weight data exist in the system, so no
motion can be singled out---every allowed motion is instantiated (full
branching), with multiplicity given by counting. The only available
measure on \(\mathcal{W}_n\) is the uniform counting measure

\[
\mu(A) = \frac{|A|}{|\mathcal{W}_n|}, \qquad |\mathcal{W}_n| = \binom{n}{n/2}.
\tag{2.4}
\]

\textbf{Proposition 2.1 (Uniqueness of the measure).} The uniform
counting measure (2.4) is not an assumption.

\emph{Proof.} Let \(\mu'\) be any measure on \(\mathcal{W}_n\). If
\(\mu'\neq\mu\), then there exist \(A,B\) such that
\(\mu'(A)/|A|\neq\mu'(B)/|B|\), i.e., there is a bias toward some
motions. By the third postulate of the axiom (no selection rule),
\(\mu'\) is incompatible with the axiom. \(\blacksquare\)

\textbf{Remark.} The third postulate is a minimality argument: it
establishes that no weight data exist within a single-axiom system,
rather than asserting that probability is physically impossible. If
probability were allowed, the full-branching conclusion would fail.

\begin{center}\rule{0.5\linewidth}{0.5pt}\end{center}

\section{Derived Structure}\label{derived-structure}

\subsection{Words, Graphs, and the Divergence
Dictionary}\label{words-graphs-and-the-divergence-dictionary}

Motion = one move along a directed edge of the graph \(\Gamma=(C,E)\);
the sequence of steps forms a word \(w=(e_1,\dots,e_n)\).

\textbf{Lemma 3.1 (Update rule).} Upon appending an edge \(u\to v\), the
divergence updates as \(d\mapsto d+e_v-e_u\).

\textbf{Lemma 3.2 (Zero-sum identity).} For any word \(w\),
\(\sum_{v\in C}d_w(v)\equiv 0\).

\emph{Proof.} Each step simultaneously ``enters one place and leaves
another,'' contributing \(+1-1=0\); induction on the number of steps
gives the result. \(\blacksquare\)

\textbf{Lemma 3.3 (Reachable set).} All reachable divergences form the
zero-sum hyperplane

\[
\{d_w\} = \Bigl\{x\in\mathbb Z^C : \sum_v x_v=0\Bigr\} = \operatorname{im} B =: H_Q, \qquad \dim H_Q = |C|-1,
\tag{3.1}
\]

where \(B\) is the incidence matrix of the graph.

\emph{Proof.} By Lemma 3.2, \(\{d_w\}\subseteq H_Q\). Conversely, when
\(\Gamma\) is connected, the image of the incidence matrix \(B\) spans
\(H_Q\), and \(B\) is totally unimodular, giving integer reachability
{[}3{]}. \(\blacksquare\)

\textbf{Remark.} Lemmas 3.1--3.3 reduce the zero-sum constraint from an
axiom to a theorem.

\subsection{Full Branching and Integer
Multiplicity}\label{full-branching-and-integer-multiplicity}

By the no-selection-rule postulate: no weights, probabilities, or
choices exist \(\Rightarrow\) every allowed motion is instantiated.
Multiplicity is given by counting, hence is integer. This is the
watershed between our framework and ``borrowed-route'' approaches: the
latter list full-branching generative semantics as an input; under our
axiom it is a corollary.

\subsection{Terminal States and
Reseeding}\label{terminal-states-and-reseeding}

\textbf{Proposition 3.4 (Necessity of reseeding).} If the active layer
cannot remain empty forever (second postulate of the axiom), then
finite-lifetime models must be accompanied by reseeding.

\emph{Proof.} A four-model comparison shows: finite-lifetime models
without reseeding necessarily extinguish; periodic-clearance plus
reseeding models persist. \(\blacksquare\)

\textbf{Proposition 3.5 (Geometric necessity of terminal states).}
Requiring a stable spectral dimension forces loop-length truncation;
branches that do not close enter a terminal state at their lifetime.

\emph{Proof.} In the closed-loop graph (nodes = rotation classes of
zero-sum cyclic words of period \(T\)), without truncation the average
degree grows unbounded with \(T\) (\(3.60\) at \(T=8\), \(9.50\) at
\(T=18\)), leading to non-manifold discretization and no stable spectral
dimension. \(\blacksquare\)

\textbf{Remark.} This is the only conditional derivation in the
framework: if the geometric requirement is abandoned, the terminal
postulate loses its justification.

\begin{center}\rule{0.5\linewidth}{0.5pt}\end{center}

\section{Layered Indices}\label{layered-indices}

The physical world resides in the evolution layer. The same quantity
manifests differently at different layers; layers must be aligned before
comparison.

\textbf{Definition 4.1 (Four-layer index).}

\begin{itemize}
\tightlist
\item
  \(\mathrm{L0}\): axiom layer---postulates, integer multiplicities,
  local compensatory moves;
\item
  \(\mathrm{L1}\): exact words (historical layer)---only increasing,
  irreversible;
\item
  \(\mathrm{L1}'\): global closed-class layer---recording and readout;
\item
  \(\mathrm{L2}\): evolution layer (active layer)---has its own update
  law.
\end{itemize}

\textbf{Proposition 4.2 (Stratification Precedence).} When two
conclusions conflict, the default explanation is not contradiction but
different layers; only opposites at the same layer constitute a genuine
contradiction.

\emph{Reasoning.} The layer index is part of the assertion; an assertion
omitting the layer index is incomplete. \(\blacksquare\)

\textbf{Corollary 4.3.} The ``readout'' (\(\pi\)) does not occupy a
layer index---it is a named input (§6).

\emph{Reasoning.} If inputs were treated as layers, any conflict between
two conclusions could be dissolved by ``different layers,'' which is the
opposite of the self-consistency requirement.

\begin{center}\rule{0.5\linewidth}{0.5pt}\end{center}

\section{Object Dictionary}\label{object-dictionary}

\subsection{Kinematic Objects}\label{kinematic-objects}

{\def\LTcaptype{none} % do not increment counter
\begin{longtable}[]{@{}
  >{\raggedright\arraybackslash}p{(\linewidth - 4\tabcolsep) * \real{0.3333}}
  >{\raggedright\arraybackslash}p{(\linewidth - 4\tabcolsep) * \real{0.3333}}
  >{\raggedright\arraybackslash}p{(\linewidth - 4\tabcolsep) * \real{0.3333}}@{}}
\toprule\noalign{}
\begin{minipage}[b]{\linewidth}\raggedright
Symbol
\end{minipage} & \begin{minipage}[b]{\linewidth}\raggedright
Name
\end{minipage} & \begin{minipage}[b]{\linewidth}\raggedright
Definition
\end{minipage} \\
\midrule\noalign{}
\endhead
\bottomrule\noalign{}
\endlastfoot
\(\Gamma=(C,E)\) & graph & positions and steps; finite, connected \\
\(w=(w_1,\dots,w_n)\) & word & sequence of steps \\
\(d_w\) & divergence & entering count \(-\) leaving count \\
\(H_Q\) & zero-sum hyperplane & \(\{x:\sum x_v=0\}\) \\
\(\mathcal{W}_n\) & balanced word space &
\(\{w\in\{\pm1\}^n:\sum w_k=0\}\) \\
\(s_k\) & prefix sum & \(\sum_{i\le k}w_i\), \(s_0=s_n=0\) \\
\end{longtable}
}

\subsection{Intrinsic Quantities}\label{intrinsic-quantities}

\[
\text{height}: h(w):=\max_{0\le k\le n}|s_k|, \qquad
\text{area}: a(w):=\sum_{k=1}^n s_k.
\tag{5.1}
\]

\textbf{Proposition 5.1.} Both \(h\) and \(a\) are anchored at \(s_0=0\)
(i.e., zero), not at \(n\).

\emph{Proof.} Both are given directly by the prefix-sum sequence, which
starts at \(s_0=0\). \(\blacksquare\)

\subsection{Readout Objects}\label{readout-objects}

{\def\LTcaptype{none} % do not increment counter
\begin{longtable}[]{@{}
  >{\raggedright\arraybackslash}p{(\linewidth - 4\tabcolsep) * \real{0.3333}}
  >{\raggedright\arraybackslash}p{(\linewidth - 4\tabcolsep) * \real{0.3333}}
  >{\raggedright\arraybackslash}p{(\linewidth - 4\tabcolsep) * \real{0.3333}}@{}}
\toprule\noalign{}
\begin{minipage}[b]{\linewidth}\raggedright
Symbol
\end{minipage} & \begin{minipage}[b]{\linewidth}\raggedright
Name
\end{minipage} & \begin{minipage}[b]{\linewidth}\raggedright
Definition
\end{minipage} \\
\midrule\noalign{}
\endhead
\bottomrule\noalign{}
\endlastfoot
\(\pi\) & partition (readout) & \(\pi:\mathcal{W}_n\to\{1,\dots,k\}\) \\
\(\mathcal{B}_b\) & block & \(\{w:\pi(w)=b\}\) \\
\(\omega_b\) & block weight &
\(\rvert\mathcal{B}_b\rvert/\rvert\mathcal{W}_n\rvert\) \\
\(q\) & participation index & \(\sum_b\omega_b^2\) \\
\(F_D\) & dimension-selection functional & \(\binom{D}{2}q^D\) \\
\(\lambda\) & normalized spectrum & top three block weights, descending
and normalized \\
\end{longtable}
}

\begin{center}\rule{0.5\linewidth}{0.5pt}\end{center}

\section{Parameter Account}\label{parameter-account}

\[
Z_0 + \underbrace{n}_{\text{free parameter}} + \underbrace{\pi}_{\text{structural input}} + \underbrace{\kappa}_{\text{scale input}} = \text{complete framework}.
\tag{6.1}
\]

\textbf{That is: one axiom, one free parameter, two named inputs.}

{\def\LTcaptype{none} % do not increment counter
\begin{longtable}[]{@{}
  >{\raggedright\arraybackslash}p{(\linewidth - 4\tabcolsep) * \real{0.3333}}
  >{\raggedright\arraybackslash}p{(\linewidth - 4\tabcolsep) * \real{0.3333}}
  >{\raggedright\arraybackslash}p{(\linewidth - 4\tabcolsep) * \real{0.3333}}@{}}
\toprule\noalign{}
\begin{minipage}[b]{\linewidth}\raggedright
Item
\end{minipage} & \begin{minipage}[b]{\linewidth}\raggedright
Type
\end{minipage} & \begin{minipage}[b]{\linewidth}\raggedright
What it buys
\end{minipage} \\
\midrule\noalign{}
\endhead
\bottomrule\noalign{}
\endlastfoot
\(Z_0\) (three postulates) & axiom & words/graph/steps; closure; full
branching; integer multiplicities \\
\(n\) & free parameter & terminal scale; entry point for blocking \\
\(\pi\) & named input & dimension-selection peak; contextuality
threshold; module flow type \\
\(\kappa\) & named input & scale (1 step = how many physical units) \\
\(\beta\varepsilon\) & parameter & effective ladder temperature (height
cost), affects state layer only \\
\(\Gamma\) & named input & topological carrier; specialized to path
graph \(\mathbb Z\) in this work \\
\end{longtable}
}

\textbf{Proven non-derivable items:} absolute scale, \(D\) itself (all
metric-side structures are graph-independent, so low dimensions cannot
be excluded), measure normalization, absolute amplitude.

\begin{center}\rule{0.5\linewidth}{0.5pt}\end{center}

\section{Main Derivation Chain}\label{main-derivation-chain}

\subsection{Kinematic Starting
Point}\label{kinematic-starting-point}

Specializing to the path graph:

\[
\Gamma = \text{path graph } 0,\pm1,\pm2,\dots; \quad \text{step} = \pm1; \quad \text{closed} \iff s_n=0.
\tag{7.1}
\]

\[
\mathcal{W}_n = \Bigl\{w\in\{\pm1\}^n:\sum_{k=1}^n w_k=0\Bigr\}, \qquad \rvert\mathcal{W}_n\rvert = \binom{n}{n/2}.
\tag{7.2}
\]

\textbf{Remark.} \(\mathcal{W}_n\) is a set selected by the zero-sum
constraint---no time, no probability, no space.

\subsection{Partition and Self-Identification
Criteria}\label{partition-and-self-identification-criteria}

Partition \(\pi:\mathcal{W}_n\to\{1,\dots,k\}\), block weights
\(\omega_b=\rvert\mathcal{B}_b\rvert/\rvert\mathcal{W}_n\rvert\).

\textbf{Definition 7.1 (Self-identification criterion).} Partition
labels must be recoverable from the block itself (after closure the
system has no memory; labels cannot be carried by memory). This requires
three structural criteria:

\begin{enumerate}
\def\labelenumi{\arabic{enumi}.}
\tightlist
\item
  \textbf{Locally Computable}: the primitive \(h(w)\) can be computed
  from a single word;
\item
  \textbf{Anchored at zero}: \(h(w)\) is referenced to \(s_0=0\), not to
  \(n\);
\item
  \textbf{Non-Vanishing Bulk}: \(\max_b\omega_b\) does not decay with
  \(n\).
\end{enumerate}

\textbf{Proposition 7.2.} The family of candidate partitions satisfying
these three criteria includes but is not limited to: height-level
partition (hmax), longest-run partition (maxrun), first-flip partition
(firstflip), and flip-count partition (nflip).

\emph{Proof.} Direct verification from the criteria. \(\blacksquare\)

\textbf{Remark.} Criteria 1 and 2 constrain the single-word primitive
\(h(w)\); the merging threshold \(t^*=\arg\max_h b_h\) (Eq. 7.6) is an
ensemble-level coarse-graining rule that depends on the global mode and
\(n\), not on individual words. The two are at different layers and do
not conflict.

\subsection{Height Blocks and the Reflection
Principle}\label{height-blocks-and-the-reflection-principle}

\[
b_h = \bigl\lvert\{w\in\mathcal{W}_n : h(w)=h\}\bigr\rvert = N_\le(n,h)-N_\le(n,h-1).
\tag{7.3}
\]

The reflection principle under two-sided absorbing boundaries
\(\pm(h+1)\) gives {[}4{]}:

\[
N_\le(n,h) = \sum_{m\in\mathbb Z}(-1)^m\binom{n}{n/2+m(h+1)}.
\tag{7.4}
\]

\textbf{Explicit computation for \(n=16\):}

{\def\LTcaptype{none} % do not increment counter
\begin{longtable}[]{@{}lrrrrrrrr@{}}
\toprule\noalign{}
\(h\) & 1 & 2 & 3 & 4 & 5 & 6 & 7 & 8 \\
\midrule\noalign{}
\endhead
\bottomrule\noalign{}
\endlastfoot
\(b_h\) & 256 & 4118 & 4858 & 2518 & 880 & 208 & 30 & 2 \\
\end{longtable}
}

\[
\sum_h b_h = 12870 = \binom{16}{8} \quad\checkmark
\tag{7.5}
\]

\textbf{Merging rule} (taking the modal level as threshold):

\[
t^* := \arg\max_h b_h; \quad h>t^* \text{ form independent blocks}; \quad h\le t^* \text{ merge into one block}.
\tag{7.6}
\]

For \(n=16\), \(t^*=3\), and the unnormalized block sizes are

\[
[\rvert\mathcal{B}_b\rvert] = [9232,\;2518,\;880,\;208,\;30,\;2], \qquad k=6.
\tag{7.7}
\]

\subsection{Three Quantum Gates}\label{three-quantum-gates}

Let \(\lambda=(\lambda_1,\lambda_2,\lambda_3)\) be the top three block
weights in descending order, normalized, and \(q=\sum_b\omega_b^2\).

\textbf{Definition 7.3 (Three gates).}

\begin{itemize}
\tightlist
\item
  \textbf{Gate 1 (survival window)}: \(q\in(1/2,\,3/5)\);
\item
  \textbf{Gate 2 (contextuality)}:
  \(S_{\max}=\mu_1\lambda_1+\mu_2(\lambda_2+\lambda_3)>2\), where
  \(\mu_1=\sqrt5\), \(\mu_2=(5-\sqrt5)/2\);
\item
  \textbf{Gate 3 (density)}:
  \(\operatorname{rank}\{\text{prime exponent difference vectors}\}=k-1\).
\end{itemize}

The \(\mu\) coefficients in Gate 2 are the von Neumann trace inequality
coefficients of the KCBS inequality {[}5,6{]}; the non-contextual bound
is \(2\).

\subsection{Dimension-Selection Functional and Peak
Condition}\label{dimension-selection-functional-and-peak-condition}

\[
F_D = \binom{D}{2}q^D.
\tag{7.8}
\]

Here \(\binom{D}{2}=\dim\mathfrak{so}(D-1,1)\) (the dimension of the
Lorentz algebra, i.e.~the number of independent Lorentz generators), and
\(q^D\) is the cost of requiring simultaneous record across \(D\)
directions.

\textbf{Lemma 7.4 (Peak condition).}

\[
\frac{F_D}{F_{D-1}} = \frac{D}{D-2}q \quad\Longrightarrow\quad \text{unique finite maximum at }D=4 \iff q\in(1/2,\,3/5).
\tag{7.9}
\]

\emph{Proof.} \(F_D/F_{D-1}>1\iff q>(D-2)/D\). At \(D=4\) this requires
\(q>1/2\); at \(D=5\), \(q<3/5\). Thus \(q\in(1/2,3/5)\) simultaneously
gives \(F_4>F_3\) and \(F_5<F_4\); for \(D\ge5\), \((D-2)/D\ge3/5>q\),
so \(F\) is strictly decreasing for \(D\ge5\). At the boundary:
\(F_0=F_1=0\), \(F_2=q^2\), \(F_3=3q^3\), and when \(q>1/2\),
\(F_3/F_2=3q>3/2>1\), so \(D\le3\) does not exceed \(F_4\). The maximum
is unique. \(\blacksquare\)

\subsection{Layer Distinction}\label{layer-distinction}

The three gates read different layers: the blocking layer (\(k\),
\(\omega_b\), \(q\)) is determined by pure counting and does not depend
on \(\beta\varepsilon\); the state layer (spectrum of \(\rho\),
\(S_{\max}\)) is determined by the measure
\(\omega\propto m\,e^{-\beta\varepsilon t}\) and varies monotonically
with \(\beta\varepsilon\).

\emph{Consistency check.} At \(n=16\), \(k=6\) for all
\(\beta\varepsilon\); \(S_{\max}\) increases monotonically from
\(2.006279\) (\(\beta\varepsilon=0\)) to \(2.235986\)
(\(\beta\varepsilon=2\)).

\subsection{\texorpdfstring{7.7 Explicit Solution for
\(n=16\)}{7.7 Explicit Solution for n=16}}\label{explicit-solution-for-n16}

\[
q = \frac{9232^2+2518^2+880^2+208^2+30^2+2^2}{12870^2} = \mathbf{0.55777859} \in (1/2,\,3/5) \quad\checkmark
\tag{7.10}
\]

\[
(\lambda_1,\lambda_2,\lambda_3) = (0.73095804,\;0.19936659,\;0.06967538),
\tag{7.11}
\]

\[
S_{\max} = \sqrt5\,(0.73095804)+\frac{5-\sqrt5}{2}(0.19936659+0.06967538) = \mathbf{2.00627871}>2 \quad\checkmark
\tag{7.12}
\]

\[
\operatorname{rank}=5=k-1 \quad\checkmark
\tag{7.13}
\]

\subsection{Full-Chain Census}\label{full-chain-census}

\[
Z_0 \longrightarrow \mathcal{W}_n \longrightarrow \pi \longrightarrow (n,\pi) \longrightarrow D=4.
\tag{7.14}
\]

A census over \(n=6,\dots,48\) (22 even values) \(\times\) 4
self-identification candidate partitions (88 combinations) yields:

\[
(n,\pi)\in\{(16,\texttt{hmax}),\;(20,\texttt{maxrun})\} \qquad (2/88 = 2.3\%).
\tag{7.15}
\]

\begin{center}\rule{0.5\linewidth}{0.5pt}\end{center}

\section{Summary of Proven
Results}\label{summary-of-proven-results}

\begin{enumerate}
\def\labelenumi{\arabic{enumi}.}
\tightlist
\item
  \textbf{One axiom}, three postulates each buying one thing (proven).
\item
  \textbf{Zero-sum constraint is a theorem, not an axiom} (Lemmas
  3.1--3.3, proven).
\item
  \textbf{Counting measure is unique} (Proposition 2.1, proven).
\item
  \textbf{Reseeding is forced by the requirement of persistent motion}
  (Proposition 3.4, proven).
\item
  \textbf{Terminal states are forced by geometric stability}
  (Proposition 3.5, conditionally derived).
\item
  \textbf{Self-identification criteria restrict the candidate partition
  family} (Proposition 7.2, proven).
\item
  \textbf{The height-level partition passes all three quantum gates and
  peaks at \(D=4\)} (§7.7, computed).
\item
  \textbf{Peak at \(D=4\iff q\in(1/2,3/5)\)} (Lemma 7.4, proven).
\item
  \textbf{\(D\) cannot be determined from proven structures; absolute
  scale is non-derivable} (§6, proven non-derivable).
\item
  \textbf{Second route}: the cycle space \(Z_1(K_4)\) and displacement
  space \(H_Q\) form an axial--polar vector duality in three dimensions
  via the Hodge star \(\star:\Lambda^2\mathbb R^3\to\mathbb R^3\),
  independently locking \(d=3\) without the \(F_D\) peak condition (see
  companion paper {[}7{]}).
\end{enumerate}

\begin{center}\rule{0.5\linewidth}{0.5pt}\end{center}

\section{Conditions and Open
Items}\label{conditions-and-open-items}

This section is of equal importance to the main text.

{\def\LTcaptype{none} % do not increment counter
\begin{longtable}[]{@{}
  >{\raggedleft\arraybackslash}p{(\linewidth - 4\tabcolsep) * \real{0.3333}}
  >{\raggedright\arraybackslash}p{(\linewidth - 4\tabcolsep) * \real{0.3333}}
  >{\raggedright\arraybackslash}p{(\linewidth - 4\tabcolsep) * \real{0.3333}}@{}}
\toprule\noalign{}
\begin{minipage}[b]{\linewidth}\raggedleft
\#
\end{minipage} & \begin{minipage}[b]{\linewidth}\raggedright
Condition
\end{minipage} & \begin{minipage}[b]{\linewidth}\raggedright
Status
\end{minipage} \\
\midrule\noalign{}
\endhead
\bottomrule\noalign{}
\endlastfoot
1 & \textbf{Quantumness required} (Gate 2: \(S_{\max}>2\)) & Premise.
Critical bottleneck of the entire chain. Necessity unproven \\
2 & \textbf{Partition criterion non-unique} & Proven non-unique: the
three structural criteria are simultaneously satisfied by hmax and
firstflip \\
3 & Necessity of the three structural criteria & Unproven (but all are
preference-free) \\
4 & Exhaustiveness of the candidate partition family & Unproven (only 4
self-identification candidates checked) \\
5 & \textbf{Conditionality of \(D=4\)} & Not true for all \(n\): only
\(n\in\{16,20\}\) give peak \(=4\). The Hodge duality route
independently locks \(d=3\), but requires the three-dimensional
geometric realization premise \\
6 & Scan range for \(n\) & Only \(n=6,\dots,48\) \\
7 & Status of \(\kappa\) & Fixed once or reselected each round;
indistinguishable within the framework \\
8 & Arrow of time & Still open \\
9 & Continuum limit / metric / field equations & Not addressed in this
work; see companion paper {[}7{]} \\
10 & \textbf{Truncation renormalization in Gate 2} & \(S_{\max}>2\) is
obtained by taking the top three blocks and renormalizing (discarding
\(h=6,7,8\), total 240 words). Without renormalization
\(S_{\max}=1.9946<2\), failing the contextuality bound. Gate 2 is
currently conditionally justified \\
11 & \textbf{Gate 2 fixed at \(\beta\varepsilon=0\)} &
\(\beta\varepsilon\) is a free parameter, but the census uniqueness
holds at \(\beta\varepsilon=0\) (pure counting); otherwise any partition
passing Gates 1 and 3 could have \(S_{\max}\) cross 2 by taking
\(\beta\varepsilon>0\) \\
12 & \textbf{Functional form \(F_D=\binom{D}{2}q^D\)} & Taking the
number of Lorentz generators times \(q^D\) and choosing the physical
dimension by the maximum is a readout functional hypothesis on par with
\(\pi\), not a theorem derived from \(Z_0\) \\
\end{longtable}
}

\begin{center}\rule{0.5\linewidth}{0.5pt}\end{center}

\section{Discussion}\label{discussion}

\[
Z_0 \Longrightarrow \mathcal{W}_n \Longrightarrow \pi \text{ (self-identifying)} \Longrightarrow D=4.
\]

All results are conditional on quantumness; the axiom itself is
unchanged.

The framework's free content consists of exactly four items: one axiom
\(Z_0\), one free parameter \(n\), two named inputs \(\pi\) and
\(\kappa\).

The two weakest links at present are: link ② (partition non-uniqueness,
proven) and the quantumness premise of link ③ (unproven). Subsequent
work targets these two links, as well as establishing the continuum
limit and field-theoretic stage.

\textbf{Claims not made:} \(Z_0\) has been modified; the chain holds
unconditionally; \(\pi\) has been uniquely determined; \(D=4\) holds for
all \(n\); contextuality resolves the arrow of time; this work provides
a continuum limit, metric, or gravitational field equations.

\begin{center}\rule{0.5\linewidth}{0.5pt}\end{center}

\section*{References}\label{references}

{[}1{]} R. Penrose, \emph{The Road to Reality: A Complete Guide to the
Laws of the Universe}, Vintage Books, 2005.

{[}2{]} C. J. Isham, ``Structural issues in quantum gravity,'' in
\emph{Recent Perspectives in the Quantum Theory of Gravity},
arXiv:gr-qc/9510034, 1994.

{[}3{]} A. Schrijver, \emph{Theory of Linear and Integer Programming},
Wiley, 1986.

{[}4{]} W. Feller, \emph{An Introduction to Probability Theory and Its
Applications}, Vol. 1, 3rd ed., Wiley, 1968 (reflection principle and
first-passage distributions).

{[}5{]} A. A. Klyachko, M. A. Can, S. Binicioğlu, A. S. Shumovsky, ``A
simple test for hidden variables in spin-1 systems,'' \emph{Phys.
Rev.~Lett.} \textbf{101}, 020403 (2008).

{[}6{]} R. W. Spekkens, ``Evidence for the epistemic view of quantum
states: A toy theory,'' \emph{Phys. Rev.~A} \textbf{75}, 032110 (2007).

{[}7{]} Companion paper: \emph{Conditional Reduction to General
Relativity from Zero-Sum Constraints: Lovelock Uniqueness and the
Group-Theoretic Origin of Dimension} (in preparation).

\end{document}
